% ---
% id: belao-2023-04
% lang: ru
% translation_status: original
% topic: radiation
% subtopic: black-body
% difficulty: 1
% source: belao
% year: 2023
% grade: 10-11
% title: Светимость бабочки семейства Papilionidae
% event: Областная олимпиада
% ---
Известно, что крылья бабочек обладают уникальными оптическими
свойствами. Несмотря на свою хрупкость, они достаточно сложны в своей
структуре, и это тенденция, присущая всем видам из всех семейств. Так
отражательная способность крыльев бабочек из семейства Papilionidae
практически равна 0, что обусловлено их полидисперсной сотовой структурой.
Учитывая, что в среднем у бабочек из семейства Papilionidae суммарная излучающая площадь крыльев
$S = \qty{20}{см}^2$, вычислите светимость бабочки. На какой спектральный диапазон
приходится максимум излучения такой бабочки? Температуру крыльев принять
равной $t = \adeg{30}\,\text{С}$. Постоянная Вина $b = 2{,}9\cdot 10^{-3}\,\text{м}\cdot\text{К}$,
постоянная Стефана-Больцмана
$\sigma = 5{,}67\cdot 10^{-8}\,\text{Вт}\cdot\text{м}^{-2}\cdot\text{К}^{-4}$.
